% PAGES: 290-290
% Last file of chunk c2 (assigned 279-290). Section 9.5 closes with (9.82);
% section 9.6 (Value at Risk) opens and runs to the end of this file. The
% file ends cleanly at a paragraph boundary ("...is $1-C$.") with no open
% \begin{...}, no open \[, and no open non-environment example block --
% chunk c3 (assigned 291-302) does not need to close anything from me.
% Rendered and read every page 279-290 individually via CropBox at 300 dpi
% (with additional 1200 dpi spot-checks where noted in earlier files in this
% chunk); confirmed check_env_balance.py result in my final report.

Thus, the asset allocation for the minimum variance portfolio fully invested in two
assets is
\begin{align}
w_1 &= \frac{\sigma_2^2-\sigma_1\sigma_2\rho_{1,2}}{\sigma_1^2-2\sigma_1\sigma_2
\rho_{1,2}+\sigma_2^2}; \\
w_2 &= \frac{\sigma_1^2-\sigma_1\sigma_2\rho_{1,2}}{\sigma_1^2-2\sigma_1\sigma_2
\rho_{1,2}+\sigma_2^2} \;=\; 1-w_1.
\end{align}

From (9.77) and (9.79), we obtain that the standard deviation of the return of this
portfolio is
\begin{equation}
\sigma_m \;=\; \frac{1}{\sqrt{\bfone^t\Sigma_R^{-1}\bfone}} \;=\;
\frac{\sigma_1\sigma_2\sqrt{1-\rho_{1,2}^2}}{\sqrt{\sigma_1^2-2\sigma_1\sigma_2
\rho_{1,2}+\sigma_2^2}}.
\end{equation}

\section{Value at Risk (VaR). Portfolio VaR.}

The Value at Risk (VaR) of a portfolio provides one number estimating, at a given
level of confidence, how much of the portfolio value could be lost over a given time
horizon. Thus, VaR depends on two parameters:
\begin{itemize}
\item $N$, the number of days in the time horizon;
\item $C$, the confidence level.
\end{itemize}

Denote by $\VaR(N,C)$ the $N$ day $C\%$ VaR of a portfolio. Then, by definition,
$\VaR(N,C)$ is the largest value (dollar amount) such that the probability that the
portfolio loss over the next $N$ days is smaller than $\VaR(N,C)$ is $C$. In other
words, if $V(0)$ and $V(N)$ denote the current value of the portfolio and the value of
the portfolio in $N$ days, respectively, then
\begin{equation}
\VaR(N,C) \;=\; \sup\{x \text{ such that } P(V(N)-V(0)\ge -x)=C\}.
\end{equation}

\begin{examplex}
If a portfolio has a five--day 98\% VaR of \$10 million, then there is a 98\%
probability that the portfolio will not lose more that \$10 million within five days
% "more that" is printed exactly that way in the source (should read "more
% than") -- transcribed as printed, not corrected.
or, equivalently, the probability of the portfolio suffering losses bigger than \$10
million within five days is 2\%.
\end{examplex}

If the portfolio value $V(N)$, or, equivalently, the portfolio return
$\frac{V(N)-V(0)}{V(0)}$, is assumed to have a continuous probability distribution
with strictly increasing probability density function, then (9.83) simplifies to:

$\VaR(N,C)$ is the number given by
\begin{equation}
P\big(V(N)-V(0) \ge -\VaR(N,C)\big) \;=\; C.
\end{equation}

Note that (9.84) is equivalent to
\begin{equation}
P\big(V(N)-V(0) \le -\VaR(N,C)\big) \;=\; 1-C,
\end{equation}
which can be interpreted as follows: the probability that the loss of the portfolio
over the next $N$ days is greater than $\VaR(N,C)$ is $1-C$.
